
\section{Next Steps}
\label{Discussions}

This section identifies three key next steps for advancing SECDA-DSE: (i) full system integration through MCP-based automation; (ii) comprehensive evaluation of both the end-to-end system and its individual components; and (iii) large-scale validation and benchmarking across multiple AI workloads, hardware platforms, and LLM configurations. 
These directions are pursued while accounting for key challenges, including the computational cost of large-scale exploration, the dependence on high-quality hardware data for effective learning, and the risk of suboptimal design generation. Next, we discuss each of these directions and challenges in detail.


\subsection{Full system integration} 

We plan to integrate full SECDA-DSE automation by connecting the DSE Explorer component with the LLM Stack through an MCP library. This library will expose SECDA-based simulation, hardware automation, the collection of generated hardware designs, and the associated resource utilization metrics per design. 
Exposing simulation endpoints will allow the LLM Stack to iteratively generate and evaluate accelerator designs without requiring full hardware synthesis, enabling efficient exploration within the SECDA simulation-driven workflow. In this setting, the LLM Stack effectively replaces the human expert within the simulation-based design loop. 

% Additionally, exposing hardware automation capabilities will allow the LLM Stack to perform both HLS and logic synthesis when required, enabling detailed evaluation through real hardware runs and FPGA fabric resource utilization. 


\subsection{Comprehensive evaluation} 

To evaluate the effectiveness of SECDA-DSE, we aim to establish a comprehensive evaluation strategy that considers both end-to-end system performance and individual subsystem contributions. 
The DSE Explorer will be evaluated based on search efficiency and parameter space coverage, while the LLM Stack will be assessed in terms of the quality of proposed refinements, validity of generated configurations, and effectiveness of retrieval-based reasoning. 

% This modular evaluation will help isolate the contribution of each module within the overall DSE pipeline and support full-stack systematic system improvement. 


\subsection{Large-scale validation and benchmarking} 

We plan to conduct large-scale validation and benchmarking of SECDA-DSE across diverse workloads, hardware platforms, and LLM configurations. This includes generating and evaluating accelerator designs for multiple AI workloads using hardware execution flows on multiple FPGA targets. 
We will explore multiple open-source LLMs (e.g., Llama, Qwen), both pre-trained and post fine-tuning, to assess performance gains or losses over time using extended datasets and cross-model validation. This comprehensive evaluation will enable systematic benchmarking of SECDA-DSE, support refinement of the fine-tuning methodology, and provide generalizability.


% Future evaluation will also compare SECDA-DSE against established DSE baselines such as Bayesian Optimization and Genetic Algorithms. We aim to assess not only final design quality but also exploration efficiency, including the number of simulations runs required to reach high-quality and valid configurations under comparable workloads as well as device constraints.

% Despite its potential, SECDA-DSE also presents several challenges. First, even simulation-based evaluation can remain computationally expensive when exploring large accelerator design spaces. 
% Second, the quality of the DSE process depends on the consistency and representativeness of the collected hardware data points used for retrieval and adaptation. 
% Finally, although the framework constrains generation through templates and validation flows, LLM-guided reasoning may still produce suboptimal or weak designs proposals, particularly in underexplored regions of the design space. 
% Addressing these challenges will be important in future iterations of the framework along with growing the database.


\subsection{Challenges} 

Despite its potential, SECDA-DSE also presents multiple challenges: (i) even simulation-based evaluation can remain computationally expensive when exploring large accelerator design spaces; 
(ii) the quality of the DSE process depends on the consistency and representativeness of the hardware data points collected used for retrieval and adaptation; 
(iii) although the SECDA framework constrains generation through templates and validation flows, LLM-guided reasoning may still produce suboptimal designs, particularly in underexplored regions of the design space. 
Addressing these challenges will be critical in future iterations of the framework, alongside expanding the hardware design database.